% ============================================================
\chapter{General Topology Review and Categorical Language}
\label{ch:review}
% ============================================================

What is a space? The question sounds simple, almost naive, yet it forced
generations of mathematicians to rethink the foundations of geometry. In the
early twentieth century, Felix Hausdorff, Maurice Fr\'echet, and others forged
the concept of a \emph{topological space} --- a structure so general that it
encompasses familiar Euclidean spaces alongside exotic objects where our
intuition falters.

This chapter is a starting point, but not a mere rehash of an introductory
course. Our goal is to revisit the foundations with fresh eyes: those of
category theory. We will reformulate classical constructions --- continuity,
products, quotients --- in a language that reveals their deep structure. And
above all, we will abandon a stubborn prejudice: the assumption that every
``reasonable'' space should be Hausdorff. Non-separated spaces lie at the
heart of this course, and that is where things get truly interesting.

\begin{intuition}
This chapter reviews the fundamental notions of general topology, reformulated in
categorical language. The goal is not to repeat an introductory course, but to set up the
vocabulary and tools needed in subsequent chapters, emphasizing aspects often neglected:
non-separated spaces, the role of the category $\mathbf{Top}$ and its properties, and the
limitations of Hausdorff-based intuition.
\end{intuition}

% ------------------------------------------------------------
\section{Topological Spaces: Definitions and First Examples}
% ------------------------------------------------------------

Everything starts with three axioms. Three deceptively simple rules that
determine which subsets of a set $X$ deserve to be called ``open.'' It is the
choice of these open sets --- not any notion of distance --- that defines the
geometry of a space. Hausdorff himself, in his \emph{Grundz\"uge der
Mengenlehre} of 1914, understood that topology is fundamentally about
\emph{neighborhoods}, not metrics.

\begin{definition}[Topological space]
A \emph{topological space} is a pair $(X,\tau)$ where $X$ is a set and
$\tau \subseteq \mathcal{P}(X)$ satisfies:
\begin{enumerate}
  \item $\emptyset \in \tau$ and $X \in \tau$;
  \item $\tau$ is closed under arbitrary unions: if $(U_i)_{i\in I} \subseteq \tau$,
    then $\bigcup_{i\in I} U_i \in \tau$;
  \item $\tau$ is closed under finite intersections: if $U_1,\ldots,U_n \in \tau$, then
    $U_1 \cap \cdots \cap U_n \in \tau$.
\end{enumerate}
The elements of $\tau$ are called \emph{open sets}. Complements of open sets are
\emph{closed sets}.
\end{definition}

The asymmetry between unions (arbitrary) and intersections (finite only) is
not a flaw in the definition: it reflects a deep reality. If we allowed
infinite intersections of open sets, the topology would collapse on itself.
In $\R$, the intersection of all intervals $(-1/n, 1/n)$ is just the
singleton $\{0\}$, which is not open. Understanding why this asymmetry is
necessary is already understanding topology.

To build intuition, let us consider the extreme cases.

\begin{example}[Extremal topologies]
On any set $X$:
\begin{itemize}
  \item The \emph{discrete topology} $\tau = \mathcal{P}(X)$ makes every subset open.
  \item The \emph{indiscrete topology} $\tau = \{\emptyset, X\}$ has only the trivial
    open sets.
\end{itemize}
\end{example}

Between these two extremes lies a tiny but remarkably rich space that will
appear again and again throughout this course.

\begin{example}[Sierpi\'nski space]\label{ex:sierpinski}
The Sierpi\'nski space is $\mathbb{S} = \{0,1\}$ with $\tau = \{\emptyset, \{1\},
\{0,1\}\}$. It is a $T_0$-space that is not $T_1$. The singleton $\{0\}$ is closed but
not open, and $\{1\}$ is open but not closed. This space plays a fundamental role
throughout the course.
\end{example}

This two-point space is in fact the ``classifier of open sets'': we will see
in the exercises that the open sets of a space $X$ are in bijection with the
continuous maps from $X$ to $\mathbb{S}$. This is the first manifestation of
a recurring theme: small non-separated spaces encode information.

\begin{example}[Cofinite topology]
On an infinite set $X$, the \emph{cofinite topology} is
$\tau = \{\emptyset\} \cup \{U \subseteq X : X \setminus U \text{ is finite}\}$.
This is a $T_1$-space that is not $T_2$.
\end{example}

% ------------------------------------------------------------
\section{Continuous Maps and Homeomorphisms}
% ------------------------------------------------------------

In topology, distance is absent, but ``nearness'' persists through open
sets. Continuity, which in analysis requires $\varepsilon$--$\delta$
arguments, takes a purer form here: a map is continuous if it respects the
open-set structure, meaning that the preimage of every open set is open.

\begin{definition}[Continuous map]
Let $(X,\tau_X)$ and $(Y,\tau_Y)$ be topological spaces. A map $f\colon X \to Y$ is
\emph{continuous} if for every $V \in \tau_Y$, we have $f^{-1}(V) \in \tau_X$.
\end{definition}

This laconic yet elegant definition admits several equivalent
reformulations, each useful in different contexts.

\begin{proposition}[Characterizations of continuity]
Let $f\colon X \to Y$. The following are equivalent:
\begin{enumerate}
  \item $f$ is continuous.
  \item For every closed set $F$ of $Y$, $f^{-1}(F)$ is closed in $X$.
  \item For every $A \subseteq X$, $f(\overline{A}) \subseteq \overline{f(A)}$.
  \item For every $x \in X$ and every neighborhood $V$ of $f(x)$, $f^{-1}(V)$ is a
    neighborhood of $x$.
\end{enumerate}
\end{proposition}

\begin{proof}
$(1) \Leftrightarrow (2)$: Immediate by taking complements.

$(1) \Rightarrow (3)$: Let $A \subseteq X$ and $x \in \overline{A}$. Let $V$ be an open
set in $Y$ containing $f(x)$. Then $f^{-1}(V)$ is open and contains $x$, so
$f^{-1}(V) \cap A \neq \emptyset$, whence $V \cap f(A) \neq \emptyset$. Thus
$f(x) \in \overline{f(A)}$.

$(3) \Rightarrow (2)$: If $F$ is closed in $Y$ and $A = f^{-1}(F)$, then
$f(\overline{A}) \subseteq \overline{f(A)} \subseteq \overline{F} = F$, so
$\overline{A} \subseteq f^{-1}(F) = A$, showing $A$ is closed.

$(1) \Leftrightarrow (4)$: Exercise.
\end{proof}

When two spaces share the same ``topological shape,'' we say they are
homeomorphic. This is the notion of equivalence in topology: the analogue of
isometry in metric geometry, or isomorphism in algebra.

\begin{definition}[Homeomorphism]
A map $f\colon X \to Y$ is a \emph{homeomorphism} if $f$ is bijective and both $f$ and
$f^{-1}$ are continuous. We then say $X$ and $Y$ are \emph{homeomorphic}, written
$X \cong Y$.
\end{definition}

% ------------------------------------------------------------
\section{Separation Axioms}
% ------------------------------------------------------------

Here is one of the central themes of this course. The separation axioms
classify topological spaces according to their ability to ``distinguish''
points. In a standard course, one focuses on Hausdorff spaces ($T_2$) and
treats the lower axioms as curiosities. But in this course, $T_0$ spaces
that fail to be $T_1$ will be our daily companions --- because they are
precisely the spaces that arise in domain theory, algebraic geometry, and the
study of spectral spaces.

\begin{definition}[Axioms $T_0$, $T_1$, $T_2$]
Let $(X,\tau)$ be a topological space.
\begin{itemize}
  \item $X$ is \emph{$T_0$} (or Kolmogorov) if for all $x \neq y$ in $X$, there exists an
    open set containing one but not the other.
  \item $X$ is \emph{$T_1$} if for all $x \neq y$, there exists an open set containing $x$
    but not $y$ (and by symmetry, one containing $y$ but not $x$).
  \item $X$ is \emph{$T_2$} (or Hausdorff) if for all $x \neq y$, there exist disjoint
    open sets $U \ni x$ and $V \ni y$.
\end{itemize}
\end{definition}

\begin{proposition}
We have the implications $T_2 \Rightarrow T_1 \Rightarrow T_0$, and no reverse
implication holds in general.
\end{proposition}

The difference between $T_0$ and $T_1$ is subtle but fundamental. In a
$T_0$ space, one can ``distinguish'' two points, but asymmetrically: one
sees an open set the other does not, but not necessarily the reverse. The
Sierpi\'nski space is the archetypal example.

\begin{remark}
A space is $T_1$ if and only if every singleton is closed. This characterization does not
extend to $T_0$-spaces: in the Sierpi\'nski space, $\{0\}$ is closed but $\{1\}$ is not.
\end{remark}

Beyond $T_2$, the hierarchy continues with regularity and normality axioms,
which concern the separation of closed sets.

\begin{definition}[Higher separation axioms]
\begin{itemize}
  \item $X$ is \emph{regular} if for every closed set $F$ and every $x \notin F$, there
    exist disjoint open sets separating $x$ and $F$. $X$ is \emph{$T_3$} if it is regular
    and $T_1$.
  \item $X$ is \emph{normal} if for every pair of disjoint closed sets $F_1, F_2$, there
    exist disjoint open sets separating them. $X$ is \emph{$T_4$} if it is normal and $T_1$.
\end{itemize}
\end{definition}

% ------------------------------------------------------------
\section{Basic Topological Constructions}
% ------------------------------------------------------------

How do we build new spaces from existing ones? Three fundamental
constructions --- products, quotients, and subspaces --- allow us to erect
the entire edifice of topology. What is remarkable is that each of these
constructions is characterized by a \emph{universal property}, a theme we
will develop fully in the categorical section.

\begin{definition}[Product topology]
Let $(X_i,\tau_i)_{i\in I}$ be a family of topological spaces. The \emph{product topology}
on $\prod_{i\in I} X_i$ is the coarsest topology making all projections
$\pi_j\colon \prod_{i\in I} X_i \to X_j$ continuous. A base for this topology consists of
sets $\prod_{i\in I} U_i$ where $U_i \in \tau_i$ and $U_i = X_i$ for all but finitely
many indices.
\end{definition}

\begin{definition}[Quotient topology]
Let $f\colon X \to Y$ be a surjection where $X$ is a topological space. The
\emph{quotient topology} on $Y$ is
$\tau_Y = \{V \subseteq Y : f^{-1}(V) \in \tau_X\}$.
This is the finest topology making $f$ continuous.
\end{definition}

\begin{definition}[Subspace]
If $(X,\tau)$ is a topological space and $A \subseteq X$, the \emph{subspace topology}
on $A$ is $\tau_A = \{U \cap A : U \in \tau\}$.
\end{definition}

% ------------------------------------------------------------
\section{Categorical Language}
% ------------------------------------------------------------

We now arrive at the tool that will transform the way we think about
topology. Category theory, born in 1945 from the work of Samuel Eilenberg
and Saunders Mac Lane, is often described as ``the mathematics of
mathematics'': it provides a language for talking about mathematical
structures and the relationships between them, abstracting away the specific
nature of the objects under study.

Why introduce this formalism in a topology course? Because the topological
constructions we have just seen --- products, quotients, subspaces --- are
special cases of categorical \emph{limits} and \emph{colimits}. Phrasing
them this way reveals why these constructions work, and more importantly,
why certain others do not.

\begin{definition}[Category]
A \emph{category} $\mathcal{C}$ consists of:
\begin{itemize}
  \item a class $\mathrm{Ob}(\mathcal{C})$ of \emph{objects};
  \item for each pair of objects $A,B$, a set $\mathrm{Hom}_{\mathcal{C}}(A,B)$ of
    \emph{morphisms};
  \item an associative composition law
    $\circ\colon \mathrm{Hom}(B,C) \times \mathrm{Hom}(A,B) \to \mathrm{Hom}(A,C)$;
  \item for each object $A$, an identity morphism $\mathrm{id}_A \in \mathrm{Hom}(A,A)$,
    neutral for composition.
\end{itemize}
\end{definition}

Here are the categories we will use throughout this course.

\begin{example}[Fundamental topological categories]
\begin{itemize}
  \item $\mathbf{Set}$: sets and functions.
  \item $\mathbf{Top}$: topological spaces and continuous maps.
  \item $\mathbf{Top}_0$: full subcategory of $\mathbf{Top}$ consisting of $T_0$-spaces.
  \item $\mathbf{Pos}$: partially ordered sets and monotone maps.
  \item $\mathbf{Lat}$: lattices and lattice morphisms.
\end{itemize}
\end{example}

To relate categories to one another, we use \emph{functors}: ``translators''
that preserve categorical structure.

\begin{definition}[Functor]
A \emph{functor} $F\colon \mathcal{C} \to \mathcal{D}$ between categories assigns:
\begin{itemize}
  \item to each object $A$ of $\mathcal{C}$, an object $F(A)$ of $\mathcal{D}$;
  \item to each morphism $f\colon A \to B$, a morphism $F(f)\colon F(A) \to F(B)$;
\end{itemize}
such that $F(\mathrm{id}_A) = \mathrm{id}_{F(A)}$ and $F(g \circ f) = F(g) \circ F(f)$.
A functor is \emph{contravariant} if it reverses the direction of arrows.
\end{definition}

And to compare two functors, we use \emph{natural transformations} --- one
of Eilenberg and Mac Lane's deepest ideas. They famously said they invented
categories in order to define natural transformations.

\begin{definition}[Natural transformation]
Let $F, G\colon \mathcal{C} \to \mathcal{D}$ be two functors. A \emph{natural
transformation} $\eta\colon F \Rightarrow G$ is a family of morphisms
$(\eta_A\colon F(A) \to G(A))_{A \in \mathrm{Ob}(\mathcal{C})}$ such that for every
morphism $f\colon A \to B$ in $\mathcal{C}$, the following diagram commutes:
\[
\begin{tikzcd}
F(A) \arrow[r, "\eta_A"] \arrow[d, "F(f)"'] & G(A) \arrow[d, "G(f)"] \\
F(B) \arrow[r, "\eta_B"'] & G(B)
\end{tikzcd}
\]
\end{definition}

The most important concept for this course is \emph{adjunction}, which
captures the duality between ``building freely'' and ``forgetting
structure.''

\begin{definition}[Adjunction]
Let $F\colon \mathcal{C} \to \mathcal{D}$ and $G\colon \mathcal{D} \to \mathcal{C}$ be
functors. We say $F$ is \emph{left adjoint} to $G$ (written $F \dashv G$) if there
exists a natural bijection
\[
  \mathrm{Hom}_{\mathcal{D}}(F(A), B) \cong \mathrm{Hom}_{\mathcal{C}}(A, G(B))
\]
for all objects $A$ of $\mathcal{C}$ and $B$ of $\mathcal{D}$.
\end{definition}

Adjunctions are everywhere in topology. Every time we encounter a
topological construction, we will be able to understand it as a left or
right adjoint. Here is why this is so powerful.

\begin{theorem}[Universal property of adjunctions]
If $F \dashv G$, then:
\begin{enumerate}
  \item $F$ preserves colimits.
  \item $G$ preserves limits.
  \item There exist natural transformations $\eta\colon \mathrm{Id}_\mathcal{C}
    \Rightarrow GF$ (unit) and $\varepsilon\colon FG \Rightarrow \mathrm{Id}_\mathcal{D}$
    (counit) satisfying the triangle identities.
\end{enumerate}
\end{theorem}

\begin{proof}
This is a classical result. The unit $\eta_A$ corresponds to the image of
$\mathrm{id}_{F(A)}$ under the bijection
$\mathrm{Hom}_\mathcal{D}(F(A),F(A)) \cong \mathrm{Hom}_\mathcal{C}(A,GF(A))$.
The counit $\varepsilon_B$ corresponds to the image of $\mathrm{id}_{G(B)}$. Preservation
of (co)limits follows because a left (resp.\ right) adjoint commutes with colimits
(resp.\ limits), since Hom-functors convert colimits to limits.
\end{proof}

% ------------------------------------------------------------
\section{$\mathbf{Top}$ as a Category: Fundamental Properties}
% ------------------------------------------------------------

Armed with this language, what can we say about the category $\mathbf{Top}$
itself? The good news is that it is very rich: it admits all conceivable
limits and colimits. The bad news --- and this is where the trouble begins
--- is that it fails to be cartesian closed.

\begin{theorem}[$\mathbf{Top}$ is complete and cocomplete]
The category $\mathbf{Top}$ admits all small limits and all small colimits. In particular:
\begin{itemize}
  \item Products exist (product topology).
  \item Equalizers exist (subspace topology).
  \item Coproducts exist (disjoint union topology).
  \item Coequalizers exist (quotient topology).
\end{itemize}
\end{theorem}

The forgetful functor $U\colon \mathbf{Top} \to \mathbf{Set}$ --- which
``forgets'' the topology and retains only the underlying set --- is
remarkably well-endowed: it has both a left and a right adjoint, which is
relatively rare.

\begin{proposition}[The forgetful functor]
The forgetful functor $U\colon \mathbf{Top} \to \mathbf{Set}$ has both a left adjoint
(discrete topology) and a right adjoint (indiscrete topology). Denoting these by $D$
and $I$:
\[
  D \dashv U \dashv I.
\]
\end{proposition}

\begin{proof}
For any space $Y$ and set $S$, a function $f\colon S \to U(Y)$ is automatically continuous
when $S$ carries the discrete topology, since every subset of $S$ is open. Thus
$\mathrm{Hom}_\mathbf{Top}(D(S),Y) \cong \mathrm{Hom}_\mathbf{Set}(S,U(Y))$, giving
$D \dashv U$.

Similarly, any function $g\colon U(X) \to S$ is continuous when $S$ carries the indiscrete
topology, since the only preimages to check are $\emptyset$ and $S$. Hence $U \dashv I$.
\end{proof}

But here is the problem that will motivate a large part of this course:
$\mathbf{Top}$ fails to be cartesian closed.

\begin{theorem}[Is $\mathbf{Top}$ cartesian closed?]
The category $\mathbf{Top}$ is \emph{not} cartesian closed. Specifically, the
compact-open topology on $\mathrm{Hom}(X,Y)$ does not generally yield an exponential
functor adjoint to the product. However, restricting to locally compact Hausdorff
spaces gives a cartesian closed subcategory.
\end{theorem}

\begin{warning}
The failure of $\mathbf{Top}$ to be cartesian closed is a fundamental limitation that
motivates the search for better-behaved topological categories, such as convergence spaces
(Chapter~7) or compactly generated spaces.
\end{warning}

% ------------------------------------------------------------
\section{The Functor $\mathcal{O}$ and the Lattice of Opens}
% ------------------------------------------------------------

Let us shift perspective. Instead of studying a space through its
\emph{points}, let us study it through its \emph{open sets}. This seemingly
innocuous idea will lead us to locale theory (Chapter~3) and constitutes one
of the deepest bridges between topology and logic.

The set of open subsets of a topological space forms a lattice with a very
particular structure: it is a \emph{frame}.

\begin{definition}[Lattice of open sets]
For a topological space $X$, the \emph{lattice of open sets} $\mathcal{O}(X)$ is the set
of open subsets of $X$, ordered by inclusion, with operations:
\begin{itemize}
  \item $\bigvee_{i\in I} U_i = \bigcup_{i\in I} U_i$ (arbitrary joins);
  \item $U \wedge V = U \cap V$ (finite meets);
  \item $\bot = \emptyset$, $\top = X$.
\end{itemize}
This is a \emph{frame}: a complete lattice satisfying the infinite distributivity law:
\[
  U \cap \bigcup_{i\in I} V_i = \bigcup_{i\in I}(U \cap V_i).
\]
\end{definition}

What makes this construction categorically interesting is that it is
functorial: continuous maps induce frame morphisms, but in the reverse
direction.

\begin{proposition}[Functoriality of $\mathcal{O}$]
The assignment $X \mapsto \mathcal{O}(X)$ extends to a contravariant functor
$\mathcal{O}\colon \mathbf{Top} \to \mathbf{Frm}^{\mathrm{op}}$.
For a continuous map $f\colon X \to Y$,
$\mathcal{O}(f) = f^{-1}\colon \mathcal{O}(Y) \to \mathcal{O}(X)$ is a frame morphism.
\end{proposition}

\begin{proof}
The inverse image $f^{-1}$ preserves arbitrary unions and finite intersections, so it is
indeed a frame morphism. Functoriality follows from
$(g \circ f)^{-1} = f^{-1} \circ g^{-1}$.
\end{proof}

\begin{remark}
This functor lies at the heart of locale theory (Chapter~3), where the perspective is
reversed: instead of studying a space via its points, one studies it via its lattice of
open sets.
\end{remark}

% ------------------------------------------------------------
\section{Initial and Final Topologies}
% ------------------------------------------------------------

To construct new topologies, two opposing strategies present themselves. The
\emph{initial} topology is the smallest topology making certain maps
continuous: this is how products and subspaces are built. The \emph{final}
topology is the largest: it is the mechanism behind quotients and disjoint
unions. These two constructions are dual, and their universal properties
encapsulate everything one needs to know to work with them.

\begin{definition}[Initial topology]
Let $X$ be a set and $(f_i\colon X \to (Y_i,\tau_i))_{i\in I}$ a family of maps. The
\emph{initial topology} on $X$ is the coarsest topology making all $f_i$ continuous. It is
generated by the subbase $\{f_i^{-1}(U) : i \in I, U \in \tau_i\}$.
\end{definition}

\begin{definition}[Final topology]
Let $Y$ be a set and $(f_i\colon (X_i,\tau_i) \to Y)_{i\in I}$ a family of maps. The
\emph{final topology} on $Y$ is the finest topology making all $f_i$ continuous:
$\tau_Y = \{V \subseteq Y : f_i^{-1}(V) \in \tau_i \text{ for all } i\}$.
\end{definition}

The power of these constructions lies in their universal properties, which
reduce continuity verification to a component-by-component test.

\begin{theorem}[Universal property]
\begin{enumerate}
  \item \emph{Initial topology:} $g\colon Z \to X$ is continuous if and only if
    $f_i \circ g$ is continuous for all $i$.
  \item \emph{Final topology:} $g\colon Y \to Z$ is continuous if and only if
    $g \circ f_i$ is continuous for all $i$.
\end{enumerate}
\end{theorem}

\begin{proof}
(1) If $g$ is continuous, the $f_i \circ g$ are continuous as compositions of continuous
maps. Conversely, for every $U \in \tau_i$, $(f_i \circ g)^{-1}(U) = g^{-1}(f_i^{-1}(U))$
is open in $Z$. Since the $f_i^{-1}(U)$ form a subbase for the initial topology, $g$ is
continuous.

(2) If $g$ is continuous, the $g \circ f_i$ are continuous. Conversely, let $V$ be open in
$Z$. For every $i$, $f_i^{-1}(g^{-1}(V)) = (g \circ f_i)^{-1}(V)$ is open in $X_i$, so
$g^{-1}(V)$ is open in the final topology.
\end{proof}

% ------------------------------------------------------------
\section{Kolmogorov Spaces and the $T_0$ Quotient}
% ------------------------------------------------------------

Let us return to the separation axioms for a construction that beautifully
illustrates categorical thinking. Given an arbitrary topological space, can
we ``make it'' $T_0$ in the most economical way possible? The answer is the
\emph{Kolmogorov quotient}, and it is an example of \emph{adjunction} in
action.

\begin{definition}[Kolmogorov quotient]
Let $(X,\tau)$ be a topological space. Define the equivalence relation:
$x \sim y$ if and only if $\overline{\{x\}} = \overline{\{y\}}$. The \emph{Kolmogorov
quotient} is $X_0 = X/{\sim}$ with the quotient topology.
\end{definition}

The idea is simple: we identify points that the topology cannot tell apart.
The result is a $T_0$ space, and this construction is ``the best possible''
in the following precise sense.

\begin{theorem}[$T_0$ reflection]
The Kolmogorov quotient $X_0$ is a $T_0$-space, and the quotient map
$q\colon X \to X_0$ is universal: for every continuous map $f\colon X \to Y$ with $Y$ a
$T_0$-space, there exists a unique continuous map $\tilde{f}\colon X_0 \to Y$ such that
$f = \tilde{f} \circ q$. Categorically, the inclusion functor
$\mathbf{Top}_0 \hookrightarrow \mathbf{Top}$ has a left adjoint.
\end{theorem}

\begin{proof}
We show $X_0$ is $T_0$. Let $[x] \neq [y]$ in $X_0$. Then
$\overline{\{x\}} \neq \overline{\{y\}}$, so there exists an open set $U$ in $X$
containing one but not the other. Since $U$ is saturated for $\sim$ (if $z \in U$ and
$\overline{\{z\}} = \overline{\{z'\}}$, then $z' \in U$), $q(U)$ is open in $X_0$ and
separates $[x]$ from $[y]$.

For universality: if $f\colon X \to Y$ is continuous and $Y$ is $T_0$, then
$\overline{\{x\}} = \overline{\{y\}}$ implies $\overline{\{f(x)\}} = \overline{\{f(y)\}}$
(by continuity), hence $f(x) = f(y)$ since $Y$ is $T_0$. Thus $f$ factors through the
quotient.
\end{proof}

% ------------------------------------------------------------
\section{Compactness and Variants}
% ------------------------------------------------------------

Let us close this review chapter with a crucial terminological clarification.
The notion of compactness, so central in topology, is subject to a
convention disagreement between the French and Anglo-Saxon traditions that
can cause serious confusion.

\begin{definition}[Compactness]
A space $X$ is \emph{compact} if every open cover of $X$ has a finite subcover.
\end{definition}

\begin{warning}
Unlike the Bourbaki convention, we do \emph{not} assume that a compact space is
Hausdorff. This convention, standard in domain theory and algebraic topology, is
essential for this course. When we need a compact Hausdorff space, we will say so
explicitly.
\end{warning}

% ------------------------------------------------------------
\section{Exercises}
% ------------------------------------------------------------

\begin{exercise}
Show that the Sierpi\'nski space $\mathbb{S}$ is the classifier of open sets: for every
space $X$, there is a bijection between $\mathcal{O}(X)$ and the set of continuous maps
$X \to \mathbb{S}$.
\end{exercise}

\begin{exercise}
Show that in the category $\mathbf{Top}$, the equalizer of $f,g\colon X
\rightrightarrows Y$ is the subspace $\{x \in X : f(x) = g(x)\}$ with the subspace
topology.
\end{exercise}

\begin{exercise}
Verify that the functor $\mathcal{O}\colon \mathbf{Top}^{\mathrm{op}} \to \mathbf{Frm}$
is faithful if and only if one restricts to $T_0$-spaces. Provide a counterexample in the
non-$T_0$ case.
\end{exercise}

\begin{exercise}
Let $X$ be a finite ordered set. Show that the Alexandrov topology on $X$ (whose open sets
are the upward-closed subsets) is the finest topology whose specialization order is the
given order.
\end{exercise}

\begin{exercise}
Show that the forgetful functor $U\colon \mathbf{Top} \to \mathbf{Set}$ preserves neither
coproducts nor coequalizers. Give explicit examples.
\end{exercise}

\begin{exercise}[Limits in $\mathbf{Top}_0$]
Show that $\mathbf{Top}_0$ is a reflective subcategory of $\mathbf{Top}$ and deduce that
it admits all limits and colimits. Do limits in $\mathbf{Top}_0$ coincide with those in
$\mathbf{Top}$?
\end{exercise}
